% Chapter 5 -- Diodes
\chapter{Diodes}\label{ch:diodes}

Every component so far has been \emph{linear}: double the voltage and you
double the current. The diode breaks that. It conducts in one direction
and blocks in the other, which sounds like a switch but behaves like
something subtler --- a one-way valve with a pressure threshold. One
junction of silicon buys you rectification, flyback protection, voltage
references, reverse-polarity guards, and light. After this chapter you
will know which of the three diode models to reach for in a calculation,
design a half-wave rectifier and estimate its ripple, add the flyback
diode that should have been next to ch.~4's relay all along, build a
zener shunt regulator and --- just as important --- know when not to.

% from-zero:begin
\begin{rulebox}[Before you start]
Chapters~2--4. \emph{Recap:} a diode is a one-way valve for current. It is the first part you meet that does not obey a straight-line $V$--$I$ rule.
\end{rulebox}
\begin{wins}[After this chapter you can]
\begin{itemize}
\item use the three diode models and choose the one that fits the question;
\item pick a resistor for an LED;
\item build a rectifier and estimate its ripple;
\item protect a circuit from reversed polarity, ESD and inductive kick-back.
\end{itemize}
\end{wins}
\begin{principle}[How to think about this]
A diode is a gate that opens only one way, and after it opens it holds the voltage almost constant. That single idea gives LEDs, rectifiers and protection circuits. Choose the simplest model that answers your question: an ideal switch for quick estimates, a fixed drop for design.
\end{principle}
% from-zero:end

\section{The junction diode: one curve, three models}\label{sec:diode-iv}

A diode is a $p$--$n$ junction: silicon doped positive on one side,
negative on the other. Its whole story is one curve, plotted in
Figure~\ref{fig:diode-iv} from a real DC sweep.

\begin{figure}[htbp]
\centering
% fig_iv.tex -- real ngspice-47 dc sweep of the chapter's diode model
\begin{tikzpicture}
\begin{axis}[
  width=0.7\linewidth, height=5.2cm,
  xlabel={diode voltage $v_D$ (V)}, ylabel={diode current $i_D$ (mA)},
  xmin=0, xmax=1.0, ymin=0, ymax=260,
  grid=both, minor tick num=1, tick label style={font=\small},
]
\addplot[blue, thick] table[x index=0, y expr=\thisrowno{1}*1e3]
  {figures/ch05/ch05_iv.dat};
\draw[dashed, gray] (axis cs:0.7,0) -- (axis cs:0.7,250)
  node[pos=0.94, anchor=south, rotate=90, font=\small] {$\approx$0.7 V knee};
\node[font=\small, anchor=west] at (axis cs:0.35,180)
  {exponential, not a wall};
\end{axis}
\end{tikzpicture}

\caption{Diode $I$--$V$ curve, real \texttt{.dc} sweep of the silicon
model used all chapter. Notice there is no wall: conduction rises
exponentially through the ``knee'' near \qty{0.7}{\volt}. At
\qty{0.6}{\volt} it passes \qty{0.014}{\milli\ampere}; at
\qty{0.7}{\volt}, \qty{0.48}{\milli\ampere}; at \qty{0.8}{\volt},
\qty{13}{\milli\ampere}.}
\label{fig:diode-iv}
\end{figure}

Physics explains the curve with the \emph{Shockley equation}: current
grows exponentially with forward voltage,
$i=I_s\bigl(e^{v/(nV_T)}-1\bigr)$, where $V_T\approx\qty{26}{\milli\volt}$
at room temperature. We state it, never derive it --- the point of a
diode in this book is the curve's \emph{shape}, not its semiconductor
derivation. What the shape buys you: current is a steep function of
voltage, so voltage is a \emph{flat} function of current. That is why a
diode makes a usable voltage reference.

Nobody computes with the exponential. You pick one of three models by
how much precision the question needs:

\begin{figure}[htbp]
\centering
% fig_models.tex -- three diode models overlaid; real curve from ngspice
\begin{tikzpicture}
\begin{axis}[
  width=0.7\linewidth, height=5.2cm,
  xlabel={$v_D$ (V)}, ylabel={$i_D$ (mA)},
  xmin=0, xmax=1.0, ymin=0, ymax=60,
  grid=both, tick label style={font=\small},
  legend style={font=\small, at={(0.03,0.97)}, anchor=north west},
]
\addplot[blue, very thick] table[x index=0, y expr=\thisrowno{1}*1e3]
  {figures/ch05/ch05_iv.dat};
\addlegendentry{real diode (Shockley)}
\addplot[black, thick, dashed, sharp plot] coordinates {(0,0) (0,60)};
\addlegendentry{ideal: 0 V drop}
\addplot[red, thick, dashed, sharp plot] coordinates {(0.7,0) (0.7,60)};
\addlegendentry{0.7 V model}
\end{axis}
\end{tikzpicture}

\caption{The three diode models on one axis. The ideal model says zero
drop whenever it conducts; the \qty{0.7}{\volt} model says the drop is a
constant; the real curve (Shockley, simulated) sits between them and is
never wrong --- it is just useless for hand arithmetic.}
\label{fig:diode-models}
\end{figure}

\begin{definition}[The three diode models]\index{diode!models}\index{ideal diode}\index{0.7 V model}
\begin{itemize}
\item \textbf{Ideal model} --- a perfect one-way switch: conducts with
\qty{0}{\volt} drop, blocks completely below it. Use for direction-of-
flow reasoning and topology questions.
\item \textbf{Constant-\qty{0.7}{\volt} model} --- conducts at a fixed
\qty{0.7}{\volt} forward drop, blocks below it. The workhorse of hand
circuit analysis; correct to within tens of millivolts in typical
circuits.
\item \textbf{Shockley model} --- the real exponential. Use only in the
simulator, where the computer does the exponential for you.
\end{itemize}
\end{definition}

\begin{checkpoint}
An LED circuit needs \qty{5}{\milli\ampere}. Which model do you use to
size the resistor --- and which does the simulator use? (Hand work:
\qty{0.7}{\volt}-style constant drop, with the LED's real $V_f$ from the
datasheet. The simulator: Shockley, always.)
\end{checkpoint}

\section{LEDs and forward voltage}\label{sec:diode-led}

An LED is a diode that wastes its recombination energy as light instead
of heat. The physics consequence you can touch: its forward drop is set
by the photon energy, so $V_f$ depends on colour --- roughly
\qty{1.8}{\volt} red, \qty{2.2}{\volt} green/yellow, \qty{3.0}{\volt} to
\qty{3.4}{\volt} blue and white. Get the number from the datasheet, not
from the generic \qty{0.7}{\volt}.

\begin{exercise}[easy]
A red LED ($V_f=\qty{1.8}{\volt}$) runs from \qty{5}{\volt} through a
resistor at \qty{10}{\milli\ampere}. Size the resistor --- once with the
correct $V_f$, once with the \qty{0.7}{\volt} silicon guess. How far
off is the current with the wrong model?
\begin{solution}
$R=(5-1.8)/0.01=\qty{320}{\ohm}$ correct. Wrong model gives
$(5-0.7)/0.01=\qty{430}{\ohm}$; real current then
$(5-1.8)/430=\qty{7.4}{\milli\ampere}$ --- 26\% low and visibly dimmer.
Using the \qty{0.7}{\volt} model on a non-silicon $V_f$ is the LED
mistake.
\end{solution}
\end{exercise}

\section{The rectifier: AC in, lumpy DC out}\label{sec:diode-rectifier}

Put a diode in series with an AC source and a load: it passes the
positive half-cycles and blocks the negative ones. That is the
\emph{half-wave rectifier} --- the simplest possible DC supply, and the
reason wall-wart chargers exist. Add a reservoir capacitor across the
load and the gaps between half-cycles get filled in from stored charge:
the capacitor discharges through the load between peaks, giving DC with
a residual \emph{ripple}.

\begin{example}[Half-wave rectifier: predict the ripple, then simulate it]\label{ex:rectifier}
\textbf{Strategy.} The output tops out near the input peak minus one
diode drop. Between peaks the capacitor discharges through the load for
almost a full period $T$; treat the discharge as roughly linear and you
get $\Delta v\approx I\,T/C$ --- the same charge-bucket arithmetic as
ch.~4's decoupling cap. Simulate to see how honest that estimate is.

\textbf{Solution.} Source: \qty{10}{\volt} peak, \qty{50}{\hertz}
($T=\qty{20}{\milli\second}$); load \qty{1}{\kilo\ohm}; $C=\qty{100}{\micro\farad}$.
Peak output $\approx 10-0.8=\qty{9.2}{\volt}$, so load current
$\approx\qty{9.2}{\milli\ampere}$ (it droops as the voltage does).
Linear estimate:
$\Delta v\approx \frac{9.2\text{m}\cdot20\text{m}}{100\mu}\approx\qty{1.8}{\volt}$
peak-to-peak. The true discharge is exponential with
$\tau=RC=\qty{100}{\milli\second}$, so the estimate is conservative ---
as the capacitor sags the load takes less current, which slows the sag.

\begin{lstlisting}[caption={\ttfamily code/ch05\_rectifier.cir --- one
diode, one capacitor, \qty{50}{\hertz} sine. The \texttt{meas} commands
quantify the ripple over the last two cycles.},label=lst:rectifier]
* ch05_rectifier.cir -- half-wave rectifier, 10 Vpk 50 Hz, 100 uF filter
.model DSI D(is=1e-14 n=1.1 rs=0.5)
Vin 1 0 sin(0 10 50)
D1 1 2 DSI
R1 2 0 1000
C1 2 0 100u
.control
tran 20u 100m
wrdata ch05_rectifier.dat v(1) v(2)
meas tran VMAX MAX v(2) from=60m to=100m
meas tran VMIN MIN v(2) from=60m to=100m
.endc
.end
\end{lstlisting}

\begin{lstlisting}[caption={Output, verbatim from ngspice-47.}]
vmax                =  9.19186e+00 at=  6.52056e-02
vmin                =  7.68742e+00 at=  6.32056e-02
\end{lstlisting}
\end{example}

Measured: $9.19-7.69=\qty{1.50}{\volt}$ p-p ripple against the predicted
$\approx\qty{1.8}{\volt}$ --- the estimate overshoots by a fifth, exactly
as the exponential argument says it should. Swap in \qty{470}{\micro\farad}
(\texttt{code/ch05\_rectifier\_bigc.cir}) and the measured ripple drops to
\qty{0.35}{\volt}; the linear estimate predicts $1.8\cdot100/470\approx
\qty{0.39}{\volt}$. Ripple scales with $1/C$ --- the estimate is a usable
upper bound, the simulator is the truth.

\begin{figure}[htbp]
\centering
% fig_rectifier.tex -- half-wave rectifier schematic + real .tran waveforms
\begin{tikzpicture}[american, font=\small]
  \draw (0,0) node[sground]{};
  \draw (0,0) to[sV, l_=$v_{\mathrm{in}}$] (0,2)
        to[D, l=D1] (2,2) -- (4,2)
        to[C, l_={\qty{100}{\micro\farad}}, *-] (4,0) -- (0,0);
  \node[font=\small, left, align=right] at (-0.5,1.0) {10 Vpk\\50 Hz};
  \draw (4,2) to[R, l={\qty{1}{\kilo\ohm}}] (6,2) -- (6,0) -- (4,0);
  \node[circ] at (4,2) {}; \node[right, font=\small] at (4.1,2.25){$v_{\mathrm{out}}$};
\end{tikzpicture}\\[0.8em]
\begin{tikzpicture}
\begin{axis}[
  width=0.72\linewidth, height=4.6cm,
  xlabel={$t$ (ms)}, ylabel={V},
  xmin=40, xmax=100, ymin=-11, ymax=11,
  grid=both, tick label style={font=\small},
  legend style={font=\small, at={(0.03,0.30)}, anchor=west},
]
\addplot[gray] table[x index=0, y index=1] {figures/ch05/ch05_rectifier_ms.dat};
\addlegendentry{$v_{\mathrm{in}}$}
\addplot[blue, thick] table[x index=0, y index=2] {figures/ch05/ch05_rectifier_ms.dat};
\addlegendentry{$v_{\mathrm{out}}$ (C=\qty{100}{\micro\farad})}
\draw[dashed, red] (axis cs:40,9.19) -- (axis cs:100,9.19);
\draw[dashed, red] (axis cs:40,7.69) -- (axis cs:100,7.69);
\node[font=\small, anchor=west, text=red] at (axis cs:50,-4)
  {ripple $\approx$ \qty{1.5}{\volt} p-p};
\end{axis}
\end{tikzpicture}

\caption{Half-wave rectifier and its real \texttt{.tran} waveforms.
Notice the output is a decaying exponential between brief recharging
peaks, not a sine minus a piece. The diode only conducts for the short
window where $v_{\mathrm{in}}$ climbs above $v_{\mathrm{out}}+0.7$.}
\label{fig:rectifier}
\end{figure}

\begin{exercise}[medium]
Same rectifier, but the load is a \qty{10}{\milli\ampere} constant sink
and $C=\qty{220}{\micro\farad}$. Predict the p-p ripple.
\begin{solution}
$\Delta v\approx IT/C=10\text{m}\cdot20\text{m}/220\mu\approx\qty{0.9}{\volt}$.
(Constant current makes the linear estimate nearly exact --- no drooping
load to save you.)
\end{solution}
\end{exercise}

\section{The flyback diode: ch.~4's missing part}\label{sec:diode-flyback}

Chapter~4 opened a relay coil and watched the switch node jump to
\qty{1646}{\volt}. In a real circuit that spike lands on a transistor
rated for maybe \qty{40}{\volt} --- once. The fix is one diode across the
coil, \emph{freewheeling}: anode at the switched node, cathode at the
supply rail. While the coil is driven, the diode is reverse-biased and
invisible. The instant the switch opens, the coil's stored current keeps
flowing --- inductor current cannot jump --- and now it has somewhere to
go: through the diode, clamping the node at rail plus one drop.

\begin{lstlisting}[caption={\ttfamily code/ch05\_flyback\_diode.cir ---
identical to ch.~4's flyback netlist except one added line: \texttt{D1 3 1
DSI}. This diode is the difference between a circuit that works and a
dead transistor.},label=lst:flybackd]
* ch05_flyback_diode.cir -- ch4 relay coil with a freewheel diode
* WARNING: omitting this diode is what kills real transistors.
.model SW1 SW vt=0.5 vh=0.1 ron=1 roff=1e9
.model DSI D(is=1e-14 n=1.1 rs=0.5)
V1 1 0 dc 12
Rcoil 1 2 70
L1 2 3 10m IC=0
S1 3 0 10 0 SW1
Cpar 3 0 100p
D1 3 1 DSI
Vsw 10 0 pulse(1 0 1m 1n 1n 1 1)
.control
tran 500n 1.06m uic
wrdata ch05_flyback_diode.dat v(3) i(l1)
meas tran VPK MAX v(3)
.endc
.end
\end{lstlisting}

\begin{lstlisting}[caption={Output, verbatim. Same relay, same switch ---
the \qty{1646}{\volt} spike is gone.}]
vpk                 =  1.30279e+01 at=  1.00001e-03
\end{lstlisting}

\qty{13.0}{\volt}: rail plus one diode drop. The coil current decays
gently through the diode instead of detonating the voltage node.

\begin{figure}[htbp]
\centering
% fig_flybackdiode.tex -- freewheel diode placement + real clamped waveform
\begin{tikzpicture}[american, font=\small]
  \draw (0,0) to[V, l_={\qty{12}{\volt}}] (0,3) -- (1.5,3)
        to[R, l_={\qty{70}{\ohm}}, a={coil $R$}] (1.5,1.6);
  \draw (1.5,1.6) to[L, l_={\qty{10}{\milli\henry}}] (1.5,0.2) -- (1.5,0);
  \draw (1.5,0) to[cspst, l_={switch}] (3.4,0) -- (4.4,0) node[sground]{};
  \draw (0,0) -- (1.5,0);
  \draw (3.4,0) to[D, l_={$D_{\mathrm{fw}}$}] (3.4,3);
  \draw (1.5,3) -- (3.4,3) -- (4.4,3);
  \draw (4.4,3) -- (4.4,0.4);
  \node[circ] at (1.5,0){}; \node[circ] at (3.4,0){}; \node[circ] at (3.4,3){};
  \node[font=\small, align=left, anchor=west] at (4.7,1.6)
    {coil current freewheels\\through $D_{\mathrm{fw}}$ at\\switch-off};
\end{tikzpicture}\\[0.8em]
\begin{tikzpicture}
\begin{axis}[
  width=0.72\linewidth, height=4.6cm,
  xlabel={$t$ (ms)}, ylabel={$v_{\mathrm{sw}}$ (V)},
  xmin=0.95, xmax=1.06, ymin=0, ymax=15,
  grid=both, tick label style={font=\small},
  restrict y to domain=-1:20,
]
\addplot[blue, thick] table[x index=0, y index=1]
  {figures/ch05/ch05_flyback_diode_ms.dat};
\draw[dashed, red] (axis cs:0.95,13.03) -- (axis cs:1.06,13.03)
  node[pos=0.2, anchor=south, font=\small] {clamped at $\approx$\qty{13}{\volt}};
\node[font=\small, anchor=south west] at (axis cs:0.96,5.5)
  {without the diode: \qty{1646}{\volt} (ch.~4 sim)};
\end{axis}
\end{tikzpicture}

\caption{Freewheel diode placement and the real clamped waveform. Notice
the diode is \emph{backwards} --- it can never conduct while the switch
is on; it exists purely for the turn-off instant.}
\label{fig:flybackd}
\end{figure}

\begin{remark}
This is not a textbook scare story. Relay drivers, solenoid valves,
motor H-bridges and DC-DC converter switches all die this way in real
products --- the failure signature is a transistor that measures
shorted drain-to-source. Every bench engineer has a ``forgot the
flyback diode'' story; the part costs cents.
\end{remark}

\begin{exercise}[hard]
Pick a flyback diode rating for our relay: coil \qty{12}{\volt},
\qty{70}{\ohm} resistance, \qty{10}{\milli\henry}, switched by a
\qty{40}{\volt}-rated transistor.
\begin{solution}
Peak diode current = coil current at switch-off =
$12/70=\qty{171}{\milli\ampere}$. Peak reverse voltage = supply =
\qty{12}{\volt} (plus margin). Any diode rated $\ge$\qty{200}{\milli\ampere},
$\ge$\qty{25}{\volt} works --- a 1N4148 (200 mA, 75 V) is the classic
answer here; a 1N4007 if the coil were fat.
\end{solution}
\end{exercise}

\section{The zener and the shunt regulator}\label{sec:diode-zener}

A \emph{zener} is a diode manufactured to break down cleanly in
\emph{reverse}: above its rated $V_z$ it conducts hard, holding its
terminals near $V_z$. Use that flat spot as a voltage reference: feed it
through a series resistor, hang your load across the zener, and it
absorbs whatever current the load doesn't take. That is the \emph{shunt
regulator}.

\begin{definition}[Zener diode]\index{zener diode}\index{shunt regulator}
A diode run intentionally in reverse breakdown, where $v\approx V_z$
over a range of currents. The shunt regulator is a series resistor plus
a zener to ground: the resistor takes the rail-to-$V_z$ difference, the
zener takes the leftover current.
\end{definition}

\begin{example}[A \qty{3.3}{\volt} reference from \qty{9}{\volt}, sized and simulated]\label{ex:zener}
\textbf{Strategy.} Size $R_s$ so the zener still sees a few milliamps at
\emph{maximum} load (otherwise it drops out of breakdown), and so its
power is safe at \emph{zero} load (when it takes everything). Then
sweep the load current and look at what ``regulates'' really means.

\textbf{Solution.} Load: up to \qty{30}{\milli\ampere}. Give the zener
$\sim$\qty{2}{\milli\ampere} margin at full load:
$R_s=(9-3.3)/(30\text{m}+2\text{m})\approx\qty{180}{\ohm}$. Worst-case
resistor power at zero load: $(5.7)^2/180\approx\qty{0.18}{\watt}$ ---
a quarter-watt part is adequate, barely. Worst-case zener power:
$3.3\cdot30\text{m}\approx\qty{0.1}{\watt}$.

\begin{lstlisting}[caption={\ttfamily code/ch05\_zener.cir --- shunt
regulator; the \texttt{.dc} sweep drives a current source as the load
and records $v_{\mathrm{out}}$.},label=lst:zener]
* ch05_zener.cir -- 3.3 V zener shunt regulator from 9 V
.model DZ3V3 D(is=1e-14 n=1.0 rs=5 bv=3.3 ibv=0.02)
V1 1 0 dc 9
Rs 1 2 180
Dz 0 2 DZ3V3
IL 2 0 dc 0
.control
dc IL 0 0.03 0.0005
wrdata ch05_zener.dat -i(v1) v(2)
.endc
.end
\end{lstlisting}

\begin{lstlisting}[caption={Selected rows of the \texttt{.dc} sweep
(every 10th of 61): load current, total supply current, output voltage.
Notice $v_{\mathrm{out}}$ slides \qty{0.22}{\volt} end to end --- this is
regulation, not magic.}]
IL (mA)   supply I (mA)   vout (V)
0.0       30.75           3.465
10.0      31.07           3.407
20.0      31.43           3.343
30.0      31.95           3.250
\end{lstlisting}
\end{example}

\begin{figure}[htbp]
\centering
% fig_zener.tex -- zener shunt regulator schematic + real load regulation
\begin{tikzpicture}[american, font=\small]
  \draw (0,0) to[V, l_={\qty{9}{\volt}}] (0,3);
  \draw (0,3) to[R, l_={$R_s$ \qty{180}{\ohm}}] (3,3);
  \draw (3,3) -- (4.5,3);
  \draw (4.5,3) to[D, a_={$D_{\mathrm{z}}$, \qty{3.3}{\volt}}] (4.5,0);
  \draw (0,0) -- (4.5,0);
  \draw (6.5,3) to[short, i=$I_L$] (6.5,2.2);
  \draw (4.5,3) -- (6.5,3);
  \draw (6.5,2.2) node[below, font=\small]{load} -- (6.5,1.4);
  \draw (6.5,1.4) -- (6.5,0) -- (4.5,0);
  \node[circ] at (4.5,3){}; \node[circ] at (4.5,0){};
  \node[above, font=\small] at (4.5,3.1){$v_{\mathrm{out}}$};
\end{tikzpicture}\\[0.8em]
\begin{tikzpicture}
\begin{axis}[
  width=0.72\linewidth, height=4.6cm,
  xlabel={load current $I_L$ (mA)}, ylabel={$v_{\mathrm{out}}$ (V)},
  xmin=0, xmax=30, ymin=3.0, ymax=3.6,
  grid=both, tick label style={font=\small},
]
\addplot[blue, thick] table[x index=0, y index=2] {figures/ch05/ch05_zener.dat};
\draw[dashed, gray] (axis cs:0,3.3) -- (axis cs:30,3.3)
  node[pos=0.65, anchor=south, font=\small] {nominal \qty{3.3}{\volt}};
\node[font=\small, anchor=west, align=left] at (axis cs:12,3.5)
  {\qty{3.47}{\volt}$\to$\qty{3.25}{\volt} over 30 mA:\\honest, not great};
\end{axis}
\end{tikzpicture}

\caption{Zener shunt regulator and its real load-regulation curve.
Notice the output sits \emph{above} \qty{3.3}{\volt} at light load and
sags below it at full load --- a low-voltage zener's knee is soft, so
``3.3 V'' is a neighborhood, not a number.}
\label{fig:zener}
\end{figure}

Honest verdict on the shunt regulator: it is simple, it is robust, it
wastes power continuously (the resistor burns $\sim$\qty{0.2}{\watt}
always), and it regulates loosely. Reach for it when you need a cheap
\emph{reference} for milliamps of bias or a comparator threshold ---
never for powering anything that cares.

\begin{exercise}[medium]
Your load can pull \qty{20}{\milli\ampere} and you want the zener to
keep $\ge$\qty{5}{\milli\ampere} at worst. What is the largest legal
$R_s$ from \qty{9}{\volt}, and what power must it survive at zero load?
\begin{solution}
$R_{s,\max}=(9-3.3)/(20\text{m}+5\text{m})=\qty{228}{\ohm}$ ---
pick \qty{220}{\ohm}. Zero-load power:
$(5.7)^2/220=\qty{0.15}{\watt}$, so a quarter-watt part survives with
40\% margin --- use a half-watt if the chassis gets hot.
\end{solution}
\end{exercise}

\section{The diode in the wild: ESD and reverse-polarity guards}\label{sec:diode-wild}

Two more jobs you will meet diodes doing on every board you open.
\emph{ESD protection}: tiny fast diodes from a signal line to each rail
clamp electrostatic hits (a finger is a few-kilovolt capacitor) before
they reach the chip. \emph{Reverse-polarity protection}: a diode in
series with the supply so a backwards plug or reversed battery does
nothing --- the diode blocks the negative rail outright.

The series guard costs you one forward drop, always. That is the whole
trade, quantified:

\begin{example}[Reverse-polarity protection: silicon vs Schottky vs no guard]\label{ex:revpol}
\textbf{Strategy.} Price the protection: what voltage does the load
lose forwards, what does it see backwards, and what power does the
diode burn. A Schottky diode trades a lower drop for looser blocking.

\textbf{Solution.} \qty{5}{\volt} supply, \qty{100}{\ohm} load
(\qty{50}{\milli\ampere} class), series diode:

\begin{lstlisting}[caption={\ttfamily code/ch05\_reverse\_pol.cir ---
\texttt{alter} re-runs the operating point with the supply flipped;
\texttt{ch05\_reverse\_pol\_sch.cir} swaps the silicon model for
Schottky.},label=lst:revpol]
* ch05_reverse_pol.cir -- series-diode reverse-polarity guard
.model DSI D(is=1e-14 n=1.1 rs=0.5)
V1 1 0 dc 5
D1 1 2 DSI
R1 2 0 100
.control
op
print v(1) v(2) i(v1)
alter v1 dc=-5
op
print v(1) v(2) i(v1)
.endc
.end
\end{lstlisting}

\begin{lstlisting}[caption={Outputs, verbatim: silicon forward,
reversed, then the Schottky run.}]
v(1) = 5.000000e+00   v(2) = 4.152554e+00   i = -4.15255e-02
v(1) = -5.00000e+00   v(2) = -5.01000e-10   i =  5.01110e-12
v(1) = 5.000000e+00   v(2) = 4.760269e+00   i = -4.76027e-02
\end{lstlisting}

Silicon: the load gets \qty{4.15}{\volt} and the diode burns
$0.85\cdot42\text{m}\approx\qty{36}{\milli\watt}$. Reversed: the load
sees $\sim$\qty{0}{\volt} and the circuit pulls picoamps. Schottky:
\qty{4.76}{\volt} delivered, only \qty{0.24}{\volt} sacrificed and
$\sim$\qty{11}{\milli\watt} burned. The trade: Schottkys leak more in
reverse and cost slightly more --- and ``just don't plug it in wrong''
remains the cheapest option on the table.
\end{example}

\section{Safety}\label{sec:diode-safety}

\begin{rulebox}
Rectifier circuits in this chapter are simulated at \qty{10}{\volt} and
lower. Real rectifiers live next to mains --- that is out of scope for
this book, full stop. The same diode that is a gentle one-way valve at
\qty{5}{\volt} is part of a lethal circuit at \qty{230}{\volt}.
\end{rulebox}

\begin{rulebox}
An inductor with no freewheel path does not negotiate --- it makes the
voltage it needs. Never switch a relay, solenoid, or motor coil without
the diode fitted first, and never trust a ``it worked on the bench''
argument: the transistor fails on the Nth switch, not the first.
\end{rulebox}

\section{Common mistakes}\label{sec:diode-mistakes}

\begin{tripwire}[Treating the diode as a zero-volt drop]
The ideal model is for direction questions only. The \qty{0.7}{\volt}
the load never sees is also power the diode has to dissipate ---
\qty{1}{\ampere} through a series guard is \qty{0.7}{\watt} to heat.
\end{tripwire}

\begin{tripwire}[Forgetting the reverse-voltage limit]
Every diode has a peak inverse voltage rating (1N4007: \qty{1000}{\volt};
a small-signal diode may only manage \qty{75}{\volt}). Exceed it and the
diode breaks down --- sometimes usefully (a zener), usually not (it
fails shorted).
\end{tripwire}

\begin{tripwire}[Sizing LED current off the ideal model]
``The diode drops zero, so $R=V/I$'' leaves nothing for the LED's real
\qty{1.8}{\volt}--\qty{3.4}{\volt} $V_f$. Always subtract $V_f$ first ---
and get it from the datasheet, because colour sets the voltage.
\end{tripwire}

\begin{tripwire}[Believing a zener holds its rated voltage]
Low-voltage zeners have soft knees: our ``3.3 V'' part delivered
\qty{3.25}{\volt} to \qty{3.47}{\volt} across the load range. If your
circuit cares about 6\%, you need a real reference IC, not a shunt
regulator.
\end{tripwire}

\section{Summary}\label{sec:diode-summary}

\begin{center}
\begin{tabularx}{\textwidth}{l>{\raggedright\arraybackslash}X}
\toprule
\textbf{Object} & \textbf{What you need to know}\\
\midrule
Diode models & ideal (0 V) for direction; \qty{0.7}{\volt} for hand
math; Shockley exponential in the simulator \\
LED & a diode that emits; $V_f$ is colour-set,
\qty{1.8}{\volt}--\qty{3.4}{\volt} --- read the datasheet \\
Rectifier & half-wave: one diode; cap filters gaps;
$\Delta v\lesssim IT/C$ p-p ripple; more $C$, less ripple \\
Flyback diode & across the coil, backwards; clamps switch-off spike at
rail+drop --- never optional \\
Zener & reverse breakdown held at $V_z$; shunt regulator = $R_s$ +
zener; loose ($\pm$5--10\%), always dissipating \\
Guards & series diode blocks reversed supply at one drop; Schottky
halves the toll; ESD diodes clamp signal pins to the rails \\
\bottomrule
\end{tabularx}
\end{center}

\section*{Exercises}
\addcontentsline{toc}{section}{Exercises}

\begin{exercise}[easy]
An ideal-model analysis says a diode drops nothing. Your circuit feeds
\qty{9}{\volt} through one diode to a \qty{3}{\volt} zener clamp. What
does the \qty{0.7}{\volt} model say the clamp sees? When does the 0.7 V
matter?
\end{exercise}

\begin{exercise}[easy]
Green LED, $V_f=\qty{2.1}{\volt}$, from \qty{3.3}{\volt} at
\qty{8}{\milli\ampere}: size the resistor.
\end{exercise}

\begin{exercise}[easy]
Which model do you use to (a) decide whether a diode conducts at all,
(b) size a series resistor by hand, (c) check a circuit in ngspice?
\end{exercise}

\begin{exercise}[easy]
In a half-wave rectifier, which half of the AC cycle appears at the
output --- and what decides which half? Trace the current path for the
conducting half.
\end{exercise}

\begin{exercise}[medium]
A rectifier must hold ripple under \qty{0.5}{\volt} p-p at
\qty{20}{\milli\ampere}, \qty{50}{\hertz}. What is the smallest
capacitor the linear estimate allows?
\end{exercise}

\begin{exercise}[medium]
Why is $\Delta v\approx IT/C$ an \emph{overestimate} of the ripple for
a resistive load but nearly exact for a constant-current load?
\end{exercise}

\begin{exercise}[medium]
A relay coil: \qty{24}{\volt}, \qty{15}{\milli\henry}, \qty{120}{\ohm},
switched at \qty{10}{\hertz}. Specify a flyback diode (current and
reverse-voltage rating with margin) and estimate the energy dissipated
per switch-off.
\end{exercise}

\begin{exercise}[medium]
A \qty{4.3}{\volt}-logic board must survive reversed \qty{5}{\volt}
input. Compare series silicon diode, series Schottky, and ``no guard''
on delivered voltage and dissipation at \qty{100}{\milli\ampere}.
\end{exercise}

\begin{exercise}[hard]
An LED matrix row drives 8 parallel LEDs ($V_f=\qty{2.0}{\volt}$ each)
from \qty{5}{\volt} through one shared resistor. Why is one resistor
per LED the correct design instead? (Hint: what happens to current
sharing when one LED's $V_f$ is \qty{0.1}{\volt} low?)
\end{exercise}

\begin{exercise}[hard]
Design a zener reference for \qty{5.1}{\volt} $\pm$\qty{0.2}{\volt}
into a load of \qty{0}{\milli\ampere}--\qty{10}{\milli\ampere} from a
\qty{12}{\volt} rail: choose $R_s$, bound the zener and resistor power,
and estimate the real output range given our measured low-$V_z$ slope
($\sim$\qty{7}{\milli\volt}/mA).
\end{exercise}

\section*{Hints and answers}
\addcontentsline{toc}{section}{Hints and answers}

\begin{description}
\item[Clamp after diode.] $9-0.7=\qty{8.3}{\volt}$ across $R_s$, so the
zener sees slightly less drive current --- the drop matters when $R_s$
is tight or the rail is low; it's why low-voltage designs use Schottky.
\item[Green LED.] $R=(3.3-2.1)/0.008=\qty{150}{\ohm}$.
\item[Model picking.] (a) ideal, (b) \qty{0.7}{\volt} (or LED $V_f$),
(c) Shockley --- the simulator never uses your shortcuts.
\item[Rectifier half.] Whichever half forward-biases the diode ---
anode side positive. Flip the diode to pick the other half.
\item[Ripple cap.] $C\ge IT/\Delta v=20\text{m}\cdot20\text{m}/0.5
=\qty{800}{\micro\farad}$ --- the estimate is an upper bound on ripple,
so this is conservative; buy \qty{1000}{\micro\farad}.
\item[Overestimate.] Resistive load current sags with the voltage, so
the bucket empties slower than $I\cdot t$ assumes; a constant-current
load never sags, so the linear discharge is real.
\item[24 V relay.] $I_{\mathrm{pk}}=24/120=\qty{200}{\milli\ampere}$;
diode $\ge$\qty{250}{\milli\ampere}, $\ge$\qty{50}{\volt} (e.g.
1N4001). Energy per switch-off:
$\frac12 LI^2=\frac12\cdot0.015\cdot0.2^2=\qty{0.3}{\milli\joule}$ ---
at \qty{10}{\hertz}, \qty{3}{\milli\watt} in the diode loop.
\item[Reversed guard.] At \qty{100}{\milli\ampere}: silicon delivers
$\approx$\qty{4.15}{\volt} and burns $\approx$\qty{85}{\milli\watt};
Schottky $\approx$\qty{4.76}{\volt}, $\approx$\qty{24}{\milli\watt};
no guard delivers \qty{5}{\volt} and dies if someone plugs it in wrong.
\item[LED matrix.] A shared resistor sets total current but the LED
with the lowest $V_f$ hogs it --- exponential $I$--$V$ means
\qty{0.1}{\volt} low doubles its share. One resistor per LED fixes
each branch's own current.
\item[5.1 V reference.] $R_s\le(12-5.1)/15\text{m}=\qty{460}{\ohm}$;
pick \qty{430}{\ohm}. Resistor power $(6.9)^2/430\approx\qty{0.11}{\watt}$;
zener worst case $5.1\cdot16\text{m}\approx\qty{0.08}{\watt}$.
At $\sim$\qty{7}{\milli\volt}/mA the output stays inside
$\pm$\qty{0.2}{\volt} up to $\sim$28 mA of load swing --- yours has 10.
Feasible.
\end{description}

\section*{Notes}
\addcontentsline{toc}{section}{Notes}

\chapternotes{%
The three-models presentation and the Shockley-stated-not-derived choice
follow the layered-model approach of the 6.002 nonlinear lectures
\cite{ocw6002} and All About Circuits vol.~III's diode chapter
\cite{aacvol3}. The flyback-diode-as-required-instrument framing is
standard bench practice; the simulated spike and clamp are real ngspice
runs of the relay from ch.~4. Zener knee softness at low $V_z$ is a
documented device property, which is why the chapter's regulator verdict
is honest rather than enthusiastic \cite{aacvol3}.}

\begin{wins}
\begin{itemize}
\item Read a diode $I$--$V$ curve and pick the right model for the job.
\item Predict and verify rectifier ripple; know why $IT/C$ overestimates.
\item Add the flyback diode to any inductive switch --- and justify its rating.
\item Build and honestly evaluate a zener shunt reference.
\item Price reverse-polarity protection in volts and milliwatts.
\end{itemize}
\end{wins}
